% arara: clean: { extensions: ['aux', 'bbl', 'blg', 'brf', 'log', 'out', 'pdf'] }
% arara: lualatex
% arara: bibtex
% arara: lualatex
% arara: lualatex
% arara: clean: { extensions: ['aux', 'bbl', 'blg', 'brf', 'log', 'out'] }
% !TeX document-id = {4a66a092-a797-425f-b588-c0b0f6bcb4d3}
% !TeX TXS-program:compile = txs:///lualatex | txs:///bibtex | txs:///lualatex | txs:///lualatex | txs:///view
%------------------------------------------------------------------------
%  Template    
%------------------------------------------------------------------------

% This file will be compiled by a computer program.
% Please do not change preamble (before II. COMMANDS ENTERED BY THE AUTHOR  ).

% Este archivo será procesado automaticamente por un programa.
% Por favor no altere nada en el preamble (hasta II. COMMANDS ENTERED BY THE AUTHOR  ).

\documentclass[11pt,a4paper]{article}
\usepackage[spanish,english]{babel}
\usepackage[T1]{fontenc} 
\usepackage{newtxtext} 
\usepackage{amsmath}

% I.5 Page layout
\usepackage[margin=3cm]{geometry}

% I.6 Title and Abstract customization
\renewenvironment{abstract}{%
	\section*{\abstractname}}
	{} 
\usepackage[affil-it]{authblk}

\usepackage[
  backref=page,
	colorlinks=true,
	urlcolor=blue,
	linkcolor=blue,
	anchorcolor=black,
	citecolor=blue
]{hyperref}

\newcommand{\MVAt}{{\usefont{U}{mvs}{m}{n}\symbol{`@}}}
% I.7 Title and author information
\title{Numerical solution for singular equations of Lane-Emden type by Adomian decomposition and Finite volume methods}
\author[1]{Matihus Alberth Molina Larios~\thanks{\href{mailto:matihus.molina@unmsm.edu.pe}{\texttt{matihus.molina\MVAt unmsm.edu.pe}}}}
\author[2]{Carlos Alonso Aznarán Laos~\thanks{\href{mailto:caznaranl@uni.pe}{\texttt{caznaranl\MVAt uni.pe}}}}
\affil[1]{Facultad de Ciencias Matemáticas, Universidad Nacional Mayor de San Marcos, Perú}
\affil[2]{Facultad de Ciencias, Universidad Nacional de Ingeniería, Perú}

\date{} %Please, do not change -  POR FAVOR, NO MODIFICAR - .
%--------------------------------------------------------------------------
%  II. COMMANDS ENTERED BY THE AUTHOR        
%--------------------------------------------------------------------------

% Please, include ONLY the COMMANDS that will be used in the abstract, if necessary.

% Por favor, incluya SOLO los COMANDOS que se utilizarán en el resumen, en caso sea necesario. 

\usepackage{filecontents}

\begin{filecontents}{\jobname.bib}
@article{Umesh2021,
	author  = {Umesh and Kumar, Manoj},
	title   = {Approximate solution of singular IVPs of Lane-Emden type and error estimation via advanced Adomian decomposition method},
	journal = {Journal of Applied Mathematics and Computing},
	year    = {2021},
	month   = {Jun},
	day     = {01},
	volume  = {66},
	number  = {1},
	pages   = {527-542},
	issn    = {1865-2085},
	doi     = {10.1007/s12190-020-01444-2}
}
@misc{Molina2025,
	title={Numerical solution for differential equations of Lane-Emden type by Adomian decomposition and integration methods}, 
	author={Molina Larios, Matihus Alberth and Aznarán Laos, Carlos Alonso},
	year={2025},
	url={https://imms.pe/documentos/26PosterMolinaLarios.pdf}
}
@article{Biles2002,
	title   = {A generalization of the Lane-Emden equation},
	journal = {Journal of Mathematical Analysis and Applications},
	volume  = {273},
	number  = {2},
	pages   = {654-666},
	year    = {2002},
	issn    = {0022-247X},
	doi     = {10.1016/S0022-247X(02)00296-2},
	author  = {Daniel C. Biles and Mark P. Robinson and John S. Spraker}
}
\end{filecontents}
\bibliographystyle{IEEEtran}
%--------------------------------------------------------------------------
%  III. DOCUMENT       
%--------------------------------------------------------------------------                                     
\pagenumbering{gobble}

\begin{document}

%The main language of the text is English. If the abstract will be written in Spanish, please uncomment the following command:
% El idioma principal del texto es English. Si el resumen será escrito en español, por favor, descomentar el siguiente comando:
%---------
%\selectlanguage{spanish}
%---------

\begin{flushleft}
	\bfseries\footnotesize VIII International Meeting in Mathematical Sciences\\
	Arequipa, Peru\\
	December 14--16, 2026
\end{flushleft}

{\let\newpage\relax\maketitle}

\begin{abstract}
	The aim of this talk is to present a well-suited approximation of a
	internal structure of a self-gravitating polytropic body model which
	consists of a type Lane-Emden equation~\cite{Umesh2021}, a second
	order nonlinear ODE.
	First, we will discuss an efficient algorithm called Adomian
	decomposition method~\cite{Molina2025} which server as an
	approximation in terms of generalized Taylor series for the
	Lane-Emden equation.
	Next, we follow the methodology given in~\cite{Biles2002} to
	guarantee the convergence to the solution.
	Finally, we carry out several numerical experiments whose results
	are in accordance with the theoretical analysis.
\end{abstract}

\bibliography{\jobname}

\end{document}
En este trabajo presentamos el problema de lane emden generalizado (\cite{les}, pero no homogéneo), que cuando m=0, e(x)
se convierte en el modelo de lame (tomamos lo que hemos escrito)
Viene el problema de valor de lane emden generalizado.
Extendemos los resultados dados por \cite{Molina,Aznarán}
en el teorema de existencia y unicidad, el sistema EDO de primer orden es 3x3.
La estimación con volumenes finitos. (Opcional: comparación con redes neuronales)
Los experimentos siguen de \cite{umesh}